\documentclass[noanswers, addpoints, 12pt]{exam}
\usepackage[utf8]{inputenc}
\usepackage[french]{babel}
\usepackage[T1]{fontenc}
\usepackage{xcolor}
\usepackage{listings}
\usepackage{graphicx}
\usepackage[most]{tcolorbox}
\usepackage{array}

\newcommand\emptybox[1]{%
 \hbox{\smash{\hspace*{\dimexpr-\hoffset}%
  }\framebox(\textwidth, #1){}}}

\definecolor{codegreen}{rgb}{0,0.6,0}
\definecolor{codegray}{rgb}{0.5,0.5,0.5}
\definecolor{codepurple}{rgb}{0.58,0,0.82}
\definecolor{backcolour}{rgb}{0.95,0.95,0.92}
 
\lstdefinestyle{mystyle}{
    backgroundcolor=\color{backcolour},   
    commentstyle=\color{codegreen},
    keywordstyle=\color{magenta},
    numberstyle=\tiny\color{codegray},
    stringstyle=\color{codepurple},
    basicstyle=\footnotesize,
    breakatwhitespace=false,         
    breaklines=true,                 
    captionpos=b,                    
    keepspaces=true,                 
    numbers=left,                    
    numbersep=5pt,                  
    showspaces=false,                
    showstringspaces=false,
    showtabs=false,                  
    tabsize=2
}
\lstset{style=mystyle, language=Python}

\title{\vspace{-2cm}Examen Première NSI - Types construits}

\date{}


\begin{document}
\runningfooter{}{\footnotesize Babash - NSI Tani Malandi et Johanne Bakalara \textperiodcentered{} CC BY-SA 4.0 \textperiodcentered{} page \thepage}{}
\firstpagefooter{}{\footnotesize Babash - NSI Tani Malandi et Johanne Bakalara \textperiodcentered{} CC BY-SA 4.0 \textperiodcentered{} page \thepage}{}

\maketitle

\pagestyle{headandfoot}
\thispagestyle{headandfoot}
\firstpageheader{Nom : \\  Prénom :\\ Classe :}{}{}
\vspace{-3cm}


\section*{Consignes:}
\begin{center} 
\fbox{\parbox{\linewidth}{
\begin{center}
Le devoir est à rédiger sur feuille libre.\\
La clarté de la rédaction et les explications font partie intégrante de la notation. Pour les questions demandant du code, respectez la syntaxe Python et soyez attentifs aux détails.
\end{center}


}}
\end{center}

\vsword{Note}
\hsword{Note}

\section*{Barème :}
\begin{center}
  \par\bigskip
  \gradetable[h]
\end{center}
\begin{questions}

\section{Comprendre les types construits en Python}

\question[1] Associer correctement les collections à leur description :
\begin{table}[h]
    \centering
    \begin{tabular}{p{3cm} p{8cm}}
        \textbf{Collection} & \textbf{Description} \\ \hline
        list & tableau mutable d’associations de paires de clé et valeur \\
        dict & collection ordonnée et non mutable d’éléments de tout type \\
        str & collection ordonnée mutable d’éléments de tout type \\
    \end{tabular}
    \label{tab:collections_descriptions}
\end{table}

\question[1] Associer correctement les collections à leurs caractéristiques. Attention, une même collection peut correspondre à plusieurs caractéristiques et inversement :
\begin{table}[h]
    \centering
    \begin{tabular}{p{3cm} p{8cm}}
        \textbf{Collection} & \textbf{Description} \\ \hline
        list & ne peut pas être modifié \\
        dict & permet un accès par index \\
        str & permet un accès par clé \\
        permet la suppression d'un élément \\
    \end{tabular}
    \label{tab:collections_caracteristiques}
\end{table}

\question[1] L’entreprise MayPy Express, spécialisée dans le transport routier de marchandises, souhaite développer un logiciel pour mieux gérer ses opérations : chauffeurs, trajets, véhicules, livraisons...
Vous faites partie de l’équipe en charge du développement.
Pour chaque situation, choisis la structure de données la plus adaptée : str, tuple, list, ou dict.
\begin{parts}
    \part Chaque matin, la secrétaire doit enregistrer les chauffeurs disponibles pour affecter des trajets. Des chauffeurs peuvent être ajoutés ou retirés au cours de la journée (ajouts ou retraits).\\
    Quelle structure de données est la plus adaptée ? Justifie brièvement.
    \part Pour optimiser les trajets, le logiciel stocke les coordonnées GPS fixes de chaque entrepôt sous forme (latitude, longitude).\\
    Quelle structure de données est la plus adaptée ? Justifie brièvement.
    \part Chaque camion a un identifiant unique (plaque d'immatriculation), et on souhaite enregistrer le nombre de livraisons effectuées par chacun des camions.\\
    Quelle structure de données est la plus adaptée ? Justifie brièvement.
    \part Chaque colis a un code de suivi unique (ex. : "XP2025AB12"). Le logiciel doit pouvoir analyser ce code caractère par caractère (par exemple, pour vérifier la structure du code).\\
    Quelle structure de données est la plus adaptée ? Justifie brièvement.
\end{parts}

\section{Identifier le type d'une variable}

\question[4] Dans le code suivant, donner le type de chacune des variables :
\begin{lstlisting}[language=Python]
message = "Bienvenue à MayPy Express"
distances = [120, 340, 85, 210]
jours_repos = ("dimanche", "lundi")
temperatures = {
    "Lyon": 18,
    "Paris": 20,
    "Toulouse": 22
}
for i in range(len(distances)):
    print(f"Trajet {i + 1} : {distances[i]} km")
for temp in temperatures.values():
    print(f"Température relevée : {temp} °C")

\end{lstlisting}
\begin{parts}
    \part Donner le type de chaucne des  variable :
    \begin{itemize}
        \item message :
        \item distances :
        \item jours\_repos :
        \item temperatures :
        \item i :
        \item temp :
    \end{itemize}
    \part Quel type de parcours est utilise à la ligne 9 ?
    \part Quel type de parcours est utilise à la ligne 11 ?
    
\end{parts}

\section{Utiliser les dictionnaires}

\question[6] L’entreprise MayPy Express gère une flotte de camions. Chaque camion est identifié par une plaque d’immatriculation, et on connaît sa capacité maximale en tonnes.

On a commencé à coder une partie du système avec un dictionnaire :
\begin{lstlisting}[language=Python]
camions = {
    "AB-123-CD": 12,
    "EF-456-GH": 7,
    "IJ-789-KL": 16
}
\end{lstlisting}
\begin{parts}
    \part Quel est le code à écrire pour obtenir la capacité du camion "EF-456-GH" ?
    \part Un nouveau camion arrive : "MN-321-OP" avec une capacité de 10 tonnes.
    Quel code faut-il écrire pour l’ajouter ?
    \part Que se passe-t-il si on exécute ce code :
    \begin{lstlisting}[language=Python]
print(camions["ZZ-999-ZZ"])\end{lstlisting}
    \part On veut afficher toutes les plaques d’immatriculation des camions.
    Écrire le code à utiliser pour cela ?
    \part Quel est le code à utiliser pour afficher chaque plaque et sa capacité, une par ligne ?
    (Exemple attendu : Camion AB-123-CD : 12 tonnes)
    \part Que va afficher ce code ?
    \begin{lstlisting}[language=Python]
for plaque, capacite in camions.items():
    if capacite >= 10:
        print(plaque)\end{lstlisting}
\end{parts}

\section{Carré magique : manipulations avancées de listes}
Un carré magique est un tableau carré (2 dimensions) contenant tous les entiers de 1 à $n^2$ une seule fois, tel que :
\begin{itemize}
    \item il contient tous les nombres entiers de 1 à $n^2$
    \item les sommes des nombres de chaque ligne, de chaque colonne et de chacune des deux diagonales sont égales.
\end{itemize}
On modélise un carré par une liste de listes de nombres. Exemples :
    \begin{lstlisting}[language=Python]
carre3 = [[2, 7, 6],
          [9, 5, 1],
          [4, 3, 8]]

carre4 = [[ 4,  5, 11, 14],
          [15, 10,  3,  1],
          [ 6,  8, 13, 12],
          [ 9, 16,  2,  7]]\end{lstlisting}
          
\question[4] Essayons de comprendre la notion de carré magique.
\begin{parts}
    \item Que vaut $carre4[2][1]$ ?
    \item Prouvez que $carre3$ est magique pour cela vérifie que chaque ligne, colonne et diagonale a la même somme.
    \item Pourquoi $carre4$ n'est-il pas magique ?
    \item Donnez un carré magique de cote 4.
\end{parts}

On propose le code suivant :
\begin{lstlisting}[language=Python]
def somme_ligne(carre: list, n: int) -> int:
    """
    carre un tableau carré de nombres
    n est un nombre entier
    """
    somme = 0
    for nombre in carre[n]:
        somme = somme + nombre
    return somme
\end{lstlisting}
\question[3] Nous cherchons maintenant à automatiser la détection d'un carré magique.
\begin{parts}
    \part Que vaut somme\_ligne(carre4, 2) ? À quoi sert cette fonction ?
    \part Écrire le code d’une fonction qui prend un carré en paramètre et qui vérifie que les sommes de nombres de chaque ligne sont égales.
    Vous pourrez utiliser la fonction somme\_ligne.
    \part Proposer le code d’une fonction somme\_colonne(carre, n) qui prend un carré en paramètre ainsi que le numéro d’une colonne, et qui renvoie la somme des nombres de cette colonne.
\end{parts}


\question[2] [Bonus] Écrivez une fonction est\_magique qui prend un carré en paramètre et renvoie True si le carré est magique et False sinon.

\end{questions}

\end{document}

